\documentclass{beamer}
\usetheme{Madrid}
\usecolortheme{default}
\setbeamertemplate{footline}[frame number]

\usepackage[T1]{fontenc}
\usepackage[utf8]{inputenc}
\usepackage{amsmath,amssymb}
\usepackage{booktabs}
\usepackage{graphicx}
\usepackage{mathtools}
\usepackage{array}

\title{Spring 2026 Natural Language Processing (CSCI-6515)}
\subtitle{Project 4: Sentiment Analysis using BERT and Reading Comprehension with BiDAF + BERT}
\author{
\textbf{Kamal Ahmadov} \\ \small kahmadov24700@ada.edu.az; kamal.ahmadov@gwu.edu \\ \and
\textbf{Rufat Guliyev} \\ \small rguliyev24988@ada.edu.az; rufat.guliyev@gwu.edu
}
\institute{ADA University \& George Washington University}
\date{April 8, 2026}

\newcommand{\EM}{\operatorname{EM}}
\newcommand{\Fone}{\operatorname{F1}}

\begin{document}

\begin{frame}
  \titlepage
\end{frame}

\begin{frame}{Outline}
\tableofcontents
\end{frame}

\section{Project Scope}

\begin{frame}{Motivation and Scope}
\begin{itemize}
  \item Project 4 combines two transformer-related tasks:
  \begin{itemize}
    \item \textbf{Task 1:} theoretical analysis of a fine-tuned BERT sentiment model
    \item \textbf{Task 2:} implemented QA system with BiDAF and frozen BERT
  \end{itemize}
  \item The shared question is the same:
  \begin{itemize}
    \item how useful are contextual transformer representations for downstream NLP tasks?
  \end{itemize}
  \item Deliverables in this repository:
  \begin{itemize}
    \item Task 1 report section
    \item Task 2 PyTorch implementation + artifacts
    \item Streamlit UI for results
  \end{itemize}
\end{itemize}
\end{frame}

\section{Task 1: BERT Sentiment Analysis}

\begin{frame}{Task 1: Selected Open-Source Model}
\small
\textbf{Model analyzed}
\begin{itemize}
  \item \texttt{nlptown/bert-base-multilingual-uncased-sentiment}
  \item Base encoder: \texttt{bert-base-multilingual-uncased}
  \item Fine-tuned for 5-class product-review sentiment
\end{itemize}

\vspace{0.4em}
\textbf{Why this model?}
\begin{itemize}
  \item Open-source and well documented
  \item More relevant to multilingual transfer than an English-only sentiment model
  \item Better fit for the Azerbaijani applicability question
\end{itemize}

\vspace{0.4em}
\textbf{Fine-tuning languages}
\begin{itemize}
  \item English, Dutch, German, French, Spanish, Italian
\end{itemize}
\end{frame}

\begin{frame}{Task 1: Inputs, Outputs, and Light Math}
\small
\textbf{Input pipeline}
\begin{itemize}
  \item Review text $\rightarrow$ WordPiece tokenizer
  \item Lowercasing because the model is \texttt{uncased}
  \item Output tensors: \texttt{input\_ids}, \texttt{attention\_mask}, optional \texttt{token\_type\_ids}
\end{itemize}

\vspace{0.4em}
\textbf{Classifier output}
\[
p(y \mid x)=\operatorname{softmax}(W h_{\texttt{[CLS]}} + b)
\]

\begin{itemize}
  \item 5 output classes: 1, 2, 3, 4, 5 stars
  \item Max input size: 512 subword positions
\end{itemize}

\vspace{0.4em}
\textbf{Case sensitivity}
\begin{itemize}
  \item Not case sensitive in practice
  \item Useful for noisy reviews, but loses emphasis from all-caps text
\end{itemize}
\end{frame}

\begin{frame}{Task 1: Can It Work for Azerbaijani?}
\small
\textbf{Short answer: yes as a baseline, no as a final solution}

\vspace{0.4em}
\textbf{Why transfer is possible}
\begin{itemize}
  \item multilingual BERT base encoder
  \item subword tokenization reduces OOV problems
  \item useful for morphologically rich languages
\end{itemize}

\vspace{0.4em}
\textbf{Why direct zero-shot use is risky}
\begin{itemize}
  \item no Azerbaijani fine-tuning
  \item trained on six other languages only
  \item Azerbaijani is agglutinative, so sentiment-bearing forms appear with many suffixes
\end{itemize}

\vspace{0.4em}
\textbf{Takeaway}
\begin{itemize}
  \item Good starting point for transfer learning
  \item Reliable Azerbaijani sentiment analysis would still require local fine-tuning
\end{itemize}
\end{frame}

\section{Task 2: Reading Comprehension}

\begin{frame}{Task 2 Goal and Variants}
\small
\textbf{Goal}
\begin{itemize}
  \item Given a question and a context passage, predict the answer span inside the context
\end{itemize}

\vspace{0.4em}
\textbf{Implemented variants}
\begin{itemize}
  \item \textbf{BiDAF + GloVe}: traditional embedding baseline
  \item \textbf{BiDAF + frozen BERT}: contextual embeddings from \texttt{bert-base-uncased}
\end{itemize}

\vspace{0.4em}
\textbf{Dataset}
\begin{itemize}
  \item SQuAD 1.1
  \item Word-level preprocessing with overlapping context windows
\end{itemize}

\vspace{0.4em}
\textbf{Core engineering choice}
\begin{itemize}
  \item Same BiDAF span-prediction stack for both variants
  \item Only the embedding source changes
\end{itemize}
\end{frame}

\begin{frame}{Task 2 Architecture (with light math)}
\small
\textbf{BiDAF pipeline}
\begin{itemize}
  \item Embeddings
  \item 2-layer highway network
  \item Contextual BiLSTM
  \item Bidirectional attention flow
  \item Modeling layers
  \item Start/end span heads
\end{itemize}

\vspace{0.4em}
\textbf{Prediction rule}
\[
(\hat{s}, \hat{e}) = \arg\max_{i \le j} \; p_{\text{start}}(i)\, p_{\text{end}}(j)
\]

\textbf{Training loss}
\[
\mathcal{L} = \operatorname{CE}(p_{\text{start}}, s^\ast) + \operatorname{CE}(p_{\text{end}}, e^\ast)
\]

\vspace{0.4em}
\textbf{BERT integration}
\begin{itemize}
  \item tokenize words into subtokens
  \item take first-subtoken hidden state for each word
  \item project BERT vectors into the BiDAF input space
\end{itemize}
\end{frame}

\begin{frame}{Task 2 Experimental Setup}
\small
\textbf{Reported run: full model settings on a larger CPU-feasible slice}
\begin{itemize}
  \item Train examples: 8192
  \item Validation examples: 2048
  \item Train windows: 9390
  \item Validation windows: 3069
  \item Epochs: 4
  \item Batch size: 16
  \item Hidden size: 64
  \item Context window: 160 words
  \item Stride: 64 words
\end{itemize}

\vspace{0.4em}
\textbf{Evaluation}
\[
\EM = \mathbf{1}[\hat{a}=a^\ast],
\qquad
\Fone = \frac{2PR}{P+R}
\]

\begin{itemize}
  \item Best checkpoint chosen by validation F1
  \item EM and F1 computed with SQuAD-style normalization
\end{itemize}
\end{frame}

\begin{frame}{Task 2 Results}
\small
\begin{table}
\centering
\begin{tabular}{lrrr}
\toprule
Variant & Best epoch & EM & F1 \\
\midrule
BiDAF + GloVe & 3 & 0.1216 & 0.1819 \\
BiDAF + frozen BERT & 4 & 0.2642 & 0.3600 \\
\bottomrule
\end{tabular}
\end{table}

\vspace{0.4em}
\textbf{Comparison}
\begin{itemize}
  \item Frozen BERT delta EM: $+0.1426$
  \item Frozen BERT delta F1: $+0.1782$
\end{itemize}

\vspace{0.4em}
\textbf{Interpretation}
\begin{itemize}
  \item The QA pipeline works end to end and scales beyond the earlier medium preset
  \item Under this larger run, frozen BERT clearly beats the traditional word-vector baseline
  \item The BERT gain is large enough to support the integration choice
\end{itemize}
\end{frame}

\begin{frame}{Task 2 Error Analysis and Engineering Notes}
\small
\textbf{Why scores still have room to improve}
\begin{itemize}
  \item still a capped subset rather than the full SQuAD training set
  \item BERT encoder is frozen, so it cannot adapt to the QA objective
  \item exact match is very strict
\end{itemize}

\vspace{0.4em}
\textbf{Important debugging issue we found}
\begin{itemize}
  \item some contexts contained standalone Unicode formatting, combining-mark, or replacement-character tokens
  \item the word tokenizer kept them, but BERT dropped them
  \item this broke word-level BERT pooling on larger reruns
\end{itemize}

\vspace{0.4em}
\textbf{Fix}
\begin{itemize}
  \item filter out tokens made only of control/format or combining-mark characters
  \item also drop standalone replacement-character tokens such as \texttt{�}
  \item after the fix, the larger compare run completes successfully
\end{itemize}
\end{frame}

\section{Extra Task: UI}

\begin{frame}{Extra Task: Streamlit UI}
\small
\textbf{Implemented UI}
\begin{itemize}
  \item \texttt{src/project4\_dashboard.py}
  \item launch with \texttt{bash scripts/run\_project4\_ui.sh}
\end{itemize}

\vspace{0.4em}
\textbf{What the UI shows}
\begin{itemize}
  \item Task 1 BERT model summary
  \item Task 2 metric boards and comparison charts
  \item prediction inspector for GloVe vs BERT
  \item artifact explorer and runbook commands
\end{itemize}

\vspace{0.4em}
\textbf{Why it matters}
\begin{itemize}
  \item easier to demonstrate during presentation
  \item easier to inspect predictions and saved outputs interactively
\end{itemize}
\end{frame}

\section{Conclusion}

\begin{frame}{Conclusion}
\small
\textbf{Task 1}
\begin{itemize}
  \item A multilingual fine-tuned BERT sentiment model has clear inputs, outputs, and limitations
  \item It is a reasonable transfer baseline for Azerbaijani, but not a complete solution
\end{itemize}

\vspace{0.4em}
\textbf{Task 2}
\begin{itemize}
  \item We implemented BiDAF in PyTorch and integrated frozen BERT embeddings
  \item We trained and evaluated on SQuAD with EM and F1
  \item In the reported larger run, frozen BERT clearly outperformed the word-vector baseline
\end{itemize}

\vspace{0.4em}
\textbf{Overall}
\begin{itemize}
  \item All major Project 4 deliverables are present in the repo:
  \begin{itemize}
    \item report
    \item code and experiment artifacts
    \item Streamlit UI
    \item presentation slides
  \end{itemize}
\end{itemize}
\end{frame}

\begin{frame}{Questions}
\centering
\Large Thank you

\vspace{0.8em}
\normalsize Questions and discussion
\end{frame}

\end{document}
